\documentclass[11pt,a4paper]{article}
\usepackage[top=1.6cm,bottom=1.4cm,left=1.3cm,right=1.3cm,headsep=8pt,footskip=16pt]{geometry}
\usepackage{enumitem}
\usepackage{amssymb}
\usepackage{array}
\usepackage{xcolor}
\usepackage{fancyhdr}
\usepackage{multicol}
\usepackage{polyglossia}
\setdefaultlanguage[calendar=gregorian, numerals=maghrib]{arabic}
\setotherlanguage{english}
\IfFontExistsTF{Amiri}{%
  \newfontfamily\arabicfont[Script=Arabic, Scale=1.08]{Amiri}
}{%
  \newfontfamily\arabicfont[Script=Arabic, Scale=1.08, Path=../fonts/, Extension=.ttf,
    UprightFont=*-Regular, BoldFont=*-Bold, ItalicFont=*-Italic,
    BoldItalicFont=*-BoldItalic]{Amiri}
}
\IfFontExistsTF{AlBattar}{%
  \newfontfamily\khatt[Script=Arabic, Scale=1.15]{AlBattar}
}{%
  \newfontfamily\khatt[Script=Arabic, Scale=1.15, Path=../fonts/,
    BoldFont=ae_AlBattar.ttf]{ae_AlBattar.ttf}
}
\renewcommand{\headrulewidth}{0.8pt}
\renewcommand{\footrulewidth}{0.8pt}
\fancyhf{}
\rhead{\khatt المدرسة العليا للأساتذة ببوزريعة}
\lhead{\khatt ملحقة بومرداس}
\lfoot{\khatt الأستاذ: د. بن عبد الرزاق يوسف}
\cfoot{\thepage/2}
\pagestyle{fancy}
\parindent 0pt
\setlength{\columnsep}{0.6cm}
\setlength{\columnseprule}{0.4pt}
\setlist[enumerate,1]{label=\textbf{\arabic*.}, leftmargin=*, itemsep=5pt, topsep=2pt}
\begin{document}

\begin{center}
{\Large\khatt\textbf{الاختبار العاشر في الذكاء الاصطناعي}}\\[2pt]
{\normalsize أسئلة متعددة الاختيارات ــ الفصول من 1 إلى 6 ــ 15 سؤالًا}
\end{center}
\vspace{-4pt}\rule{\textwidth}{0.6pt}\\
{\small\textbf{التعليمة:} لكل سؤال إجابة صحيحة واحدة فقط من بين خمسة اختيارات.
اكتب \textbf{حرف} الإجابة الصحيحة في جدول الإجابات أسفل الصفحة الثانية.
لكل سؤال نقطة واحدة (المجموع 15 نقطة).}
\vspace{-2pt}
\begin{multicols}{2}

\noindent\begin{minipage}{\linewidth}
\hangindent=0.7cm\hangafter=1\makebox[0.7cm][r]{\textbf{1.}}\textbf{ما المرحلة الأولى في تعلّم الآلة؟}
\begin{itemize}[leftmargin=0.95cm, labelsep=4pt, itemsep=0pt, topsep=2pt, parsep=0pt]
  \item[\textbf{أ)}] التدريب
  \item[\textbf{ب)}] النموذج
  \item[\textbf{ج)}] الاختبار
  \item[\textbf{د)}] البيانات
  \item[\textbf{هـ)}] النتيجة
\end{itemize}
\end{minipage}\par\vspace{7pt plus 4pt}
\noindent\begin{minipage}{\linewidth}
\hangindent=0.7cm\hangafter=1\makebox[0.7cm][r]{\textbf{2.}}\textbf{قدّم مستخدم معلومة خاطئة أخذها من أداة. من يتحمل المسؤولية؟}
\begin{itemize}[leftmargin=0.95cm, labelsep=4pt, itemsep=0pt, topsep=2pt, parsep=0pt]
  \item[\textbf{أ)}] الإنترنت
  \item[\textbf{ب)}] المستخدم
  \item[\textbf{ج)}] لا أحد
  \item[\textbf{د)}] الشركة وحدها
  \item[\textbf{هـ)}] الأداة
\end{itemize}
\end{minipage}\par\vspace{7pt plus 4pt}
\noindent\begin{minipage}{\linewidth}
\hangindent=0.7cm\hangafter=1\makebox[0.7cm][r]{\textbf{3.}}\textbf{تحويل درس «الجملة الاسمية» إلى خريطة ذهنية هو تحويل:}
\begin{itemize}[leftmargin=0.95cm, labelsep=4pt, itemsep=0pt, topsep=2pt, parsep=0pt]
  \item[\textbf{أ)}] نص $\leftarrow$ مخطط
  \item[\textbf{ب)}] صوت $\leftarrow$ نص
  \item[\textbf{ج)}] صورة $\leftarrow$ نص
  \item[\textbf{د)}] نص $\leftarrow$ صوت
  \item[\textbf{هـ)}] فيديو $\leftarrow$ صوت
\end{itemize}
\end{minipage}\par\vspace{7pt plus 4pt}
\noindent\begin{minipage}{\linewidth}
\hangindent=0.7cm\hangafter=1\makebox[0.7cm][r]{\textbf{4.}}\textbf{الشركة المطوِّرة لأداة \textenglish{Qwen} هي:}
\begin{itemize}[leftmargin=0.95cm, labelsep=4pt, itemsep=0pt, topsep=2pt, parsep=0pt]
  \item[\textbf{أ)}] \textenglish{Google}
  \item[\textbf{ب)}] \textenglish{OpenAI}
  \item[\textbf{ج)}] \textenglish{Anthropic}
  \item[\textbf{د)}] \textenglish{Microsoft}
  \item[\textbf{هـ)}] \textenglish{Alibaba}
\end{itemize}
\end{minipage}\par\vspace{7pt plus 4pt}
\noindent\begin{minipage}{\linewidth}
\hangindent=0.7cm\hangafter=1\makebox[0.7cm][r]{\textbf{5.}}\textbf{ما النموذج في مراحل تعلم الآلة؟}
\begin{itemize}[leftmargin=0.95cm, labelsep=4pt, itemsep=0pt, topsep=2pt, parsep=0pt]
  \item[\textbf{أ)}] ما ينتج عن التدريب: برنامج «تعلّم» القاعدة
  \item[\textbf{ب)}] الطلب الذي يكتبه المستخدم ليحصل على إجابة
  \item[\textbf{ج)}] مجموعة الأمثلة الأولى التي نجمعها للتدريب
  \item[\textbf{د)}] جهاز الحاسوب الذي يجري عليه التدريب كله
  \item[\textbf{هـ)}] الأمثلة الجديدة التي نختبر بها في النهاية
\end{itemize}
\end{minipage}\par\vspace{7pt plus 4pt}
\noindent\begin{minipage}{\linewidth}
\hangindent=0.7cm\hangafter=1\makebox[0.7cm][r]{\textbf{6.}}\textbf{لماذا يجب أن يكون الأمر واضحًا وكاملًا؟}
\begin{itemize}[leftmargin=0.95cm, labelsep=4pt, itemsep=0pt, topsep=2pt, parsep=0pt]
  \item[\textbf{أ)}] لأن الأداة لا تفهم إلا اللغة الإنجليزية
  \item[\textbf{ب)}] لأن الأداة لا تعرف إلا ما تكتبه لها
  \item[\textbf{ج)}] لأن الأمر الطويل أسرع في التنفيذ دائمًا
  \item[\textbf{د)}] لأن الأداة تقرأ أفكار المستخدم وحدها
  \item[\textbf{هـ)}] لأن الأوامر القصيرة ممنوعة في الأدوات
\end{itemize}
\end{minipage}\par\vspace{7pt plus 4pt}
\noindent\begin{minipage}{\linewidth}
\hangindent=0.7cm\hangafter=1\makebox[0.7cm][r]{\textbf{7.}}\textbf{ما صفات جملة الإفصاح الجيدة؟}
\begin{itemize}[leftmargin=0.95cm, labelsep=4pt, itemsep=0pt, topsep=2pt, parsep=0pt]
  \item[\textbf{أ)}] بالإنجليزية دائمًا
  \item[\textbf{ب)}] قصيرة وصادقة
  \item[\textbf{ج)}] مخفية عن القارئ
  \item[\textbf{د)}] عامة وغامضة
  \item[\textbf{هـ)}] طويلة ومبهمة
\end{itemize}
\end{minipage}\par\vspace{7pt plus 4pt}
\noindent\begin{minipage}{\linewidth}
\hangindent=0.7cm\hangafter=1\makebox[0.7cm][r]{\textbf{8.}}\textbf{ما الشفافية؟}
\begin{itemize}[leftmargin=0.95cm, labelsep=4pt, itemsep=0pt, topsep=2pt, parsep=0pt]
  \item[\textbf{أ)}] ترجمة إجابة الأداة إلى لغة القارئ
  \item[\textbf{ب)}] إخفاء طريقة استعمال الأداة عن القارئ
  \item[\textbf{ج)}] توضيح طريقة استعمال الأداة عند الحاجة
  \item[\textbf{د)}] تحمّل الأداة وحدها مسؤولية أخطائها
  \item[\textbf{هـ)}] نسخ إجابة الأداة كما هي دون تعديل
\end{itemize}
\end{minipage}\par\vspace{7pt plus 4pt}
\noindent\begin{minipage}{\linewidth}
\hangindent=0.7cm\hangafter=1\makebox[0.7cm][r]{\textbf{9.}}\textbf{ما الترتيب الصحيح لمراحل البحث عن المعلومة؟}
\begin{itemize}[leftmargin=0.95cm, labelsep=4pt, itemsep=0pt, topsep=2pt, parsep=0pt]
  \item[\textbf{أ)}] قيّم المصادر $\leftarrow$ حدّد الحاجة $\leftarrow$ ابحث $\leftarrow$ لخّص
  \item[\textbf{ب)}] ابحث $\leftarrow$ حدّد الحاجة $\leftarrow$ لخّص $\leftarrow$ قيّم المصادر
  \item[\textbf{ج)}] حدّد الحاجة $\leftarrow$ ابحث $\leftarrow$ قيّم المصادر $\leftarrow$ لخّص
  \item[\textbf{د)}] حدّد الحاجة $\leftarrow$ لخّص $\leftarrow$ ابحث $\leftarrow$ قيّم المصادر
  \item[\textbf{هـ)}] لخّص $\leftarrow$ ابحث $\leftarrow$ قيّم المصادر $\leftarrow$ حدّد الحاجة
\end{itemize}
\end{minipage}\par\vspace{7pt plus 4pt}
\noindent\begin{minipage}{\linewidth}
\hangindent=0.7cm\hangafter=1\makebox[0.7cm][r]{\textbf{10.}}\textbf{أي مما يلي من المخرجات التي تستطيع الأداة متعددة الأنماط إنتاجها حسب الدرس؟}
\begin{itemize}[leftmargin=0.95cm, labelsep=4pt, itemsep=0pt, topsep=2pt, parsep=0pt]
  \item[\textbf{أ)}] شرح مكتوب ومخطط وقراءة صوتية
  \item[\textbf{ب)}] جداول رقمية فقط دون شرح
  \item[\textbf{ج)}] أصوات وتسجيلات فقط دون نص
  \item[\textbf{د)}] صور فوتوغرافية حقيقية فقط
  \item[\textbf{هـ)}] نصوص مكتوبة فقط دون غيرها
\end{itemize}
\end{minipage}\par\vspace{7pt plus 4pt}
\noindent\begin{minipage}{\linewidth}
\hangindent=0.7cm\hangafter=1\makebox[0.7cm][r]{\textbf{11.}}\textbf{ما الذي يفسّر اختلاف إجابات الأدوات للطلب نفسه، وماذا نستنتج عمليًا؟}
\begin{itemize}[leftmargin=0.95cm, labelsep=4pt, itemsep=0pt, topsep=2pt, parsep=0pt]
  \item[\textbf{أ)}] العربية غير مدعومة، فنكتب الطلب بالإنجليزية فقط
  \item[\textbf{ب)}] الإجابات عشوائية تمامًا، فلا فائدة من الأدوات
  \item[\textbf{ج)}] إحدى الأدوات معطلة، فنعيد تشغيل الجهاز ونكرر
  \item[\textbf{د)}] لكل أداة أسلوبها، فمن المفيد المقارنة بين الأدوات
  \item[\textbf{هـ)}] الطلب خاطئ، فلا نستعمل إلا أداة واحدة بعد ذلك
\end{itemize}
\end{minipage}\par\vspace{7pt plus 4pt}
\noindent\begin{minipage}{\linewidth}
\hangindent=0.7cm\hangafter=1\makebox[0.7cm][r]{\textbf{12.}}\textbf{أداة تقترح جملًا للإعراب تدرّبت على أمثلة من بيئة المدينة فقط. ماذا قد يحدث مع تلاميذ من بيئات أخرى؟}
\begin{itemize}[leftmargin=0.95cm, labelsep=4pt, itemsep=0pt, topsep=2pt, parsep=0pt]
  \item[\textbf{أ)}] لا يحدث شيء، لأن قواعد الإعراب واحدة للجميع
  \item[\textbf{ب)}] تصبح الأمثلة أكثر تنوعًا لأن التلاميذ مختلفون
  \item[\textbf{ج)}] تتوقف الأداة عن العمل نهائيًا مع هؤلاء التلاميذ
  \item[\textbf{د)}] تتحسن الأمثلة تلقائيًا مع مرور الوقت وحدها
  \item[\textbf{هـ)}] تقترح أمثلة لا تناسبهم، وهذا شكل من الانحياز
\end{itemize}
\end{minipage}\par\vspace{7pt plus 4pt}
\noindent\begin{minipage}{\linewidth}
\hangindent=0.7cm\hangafter=1\makebox[0.7cm][r]{\textbf{13.}}\textbf{الأمر: «اشرح الاستعارة، واكتب قصيدة عن الربيع، وترجم النص إلى الفرنسية، ولخّص درس الأمس». ما الخطأ وما الحل؟}
\begin{itemize}[leftmargin=0.95cm, labelsep=4pt, itemsep=0pt, topsep=2pt, parsep=0pt]
  \item[\textbf{أ)}] غياب الدور؛ نضيف دورًا ونترك المهام
  \item[\textbf{ب)}] أمر قصير جدًا؛ نضيف إليه كلمات أكثر
  \item[\textbf{ج)}] لا خطأ فيه؛ فهو أمر نموذجي ومكتمل
  \item[\textbf{د)}] أمر غامض؛ نحذف منه كل المهام المذكورة
  \item[\textbf{هـ)}] مهام مختلطة؛ نقسمها إلى أوامر متتابعة
\end{itemize}
\end{minipage}\par\vspace{7pt plus 4pt}
\noindent\begin{minipage}{\linewidth}
\hangindent=0.7cm\hangafter=1\makebox[0.7cm][r]{\textbf{14.}}\textbf{مقال على موقع شركة تبيع برنامجًا، يؤكد أن برنامجها يضاعف مستوى التلاميذ، دون أي مرجع. أي معيارين يثيران الشك أكثر؟}
\begin{itemize}[leftmargin=0.95cm, labelsep=4pt, itemsep=0pt, topsep=2pt, parsep=0pt]
  \item[\textbf{أ)}] التاريخ (الحداثة) والكاتب (الاختصاص)
  \item[\textbf{ب)}] لا شيء، لأن الشركة معروفة ومشهورة
  \item[\textbf{ج)}] الهدف (البيع) والدقة (غياب المراجع)
  \item[\textbf{د)}] طول المقال وجودة تنسيقه وصوره
  \item[\textbf{هـ)}] الجهة وحدها لأنها ليست وزارة رسمية
\end{itemize}
\end{minipage}\par\vspace{7pt plus 4pt}
\noindent\begin{minipage}{\linewidth}
\hangindent=0.7cm\hangafter=1\makebox[0.7cm][r]{\textbf{15.}}\textbf{أي أمر أنسب لتحويل نص درس إلى مخطط مفاهيمي؟}
\begin{itemize}[leftmargin=0.95cm, labelsep=4pt, itemsep=0pt, topsep=2pt, parsep=0pt]
  \item[\textbf{أ)}] «اشرح النص المرفق شرحًا مفصلًا في عدة صفحات»
  \item[\textbf{ب)}] «ارسم لي أي شيء جميل يتعلق بهذا النص المرفق»
  \item[\textbf{ج)}] «اكتب لي مخططًا، أو جدولًا، أو قصة، كما تريد أنت»
  \item[\textbf{د)}] «حوّل النص إلى مخطط مفاهيمي يوضح العلاقة بين المفاهيم»
  \item[\textbf{هـ)}] «لخّص النص المرفق في سطر واحد فقط لا أكثر»
\end{itemize}
\end{minipage}\par\vspace{7pt plus 4pt}
\end{multicols}
\vfill
\begin{center}
{\khatt\textbf{جدول الإجابات}}\\[3pt]
\renewcommand{\arraystretch}{1.6}
\begin{tabular}{|>{\centering\arraybackslash}m{1.9cm}|*{15}{>{\centering\arraybackslash}m{0.62cm}|}}
\hline
\textbf{رقم السؤال} & \textbf{1} & \textbf{2} & \textbf{3} & \textbf{4} & \textbf{5} & \textbf{6} & \textbf{7} & \textbf{8} & \textbf{9} & \textbf{10} & \textbf{11} & \textbf{12} & \textbf{13} & \textbf{14} & \textbf{15} \\ \hline
\textbf{الإجابة} &  &  &  &  &  &  &  &  &  &  &  &  &  &  &  \\[8pt] \hline
\end{tabular}

\vspace{8pt}
\renewcommand{\arraystretch}{2.2}
\begin{tabular}{|>{\centering\arraybackslash}m{7.5cm}|>{\centering\arraybackslash}m{3.2cm}|>{\centering\arraybackslash}m{3.2cm}|}
\hline
\textbf{الاسم واللقب} & \textbf{الفوج} & \textbf{العلامة} \\ \hline
 & & \hfill /15 \\ \hline
\end{tabular}
\end{center}
\end{document}
